% Generated draft: requires mathematical and source review.
\documentclass[11pt]{article}
\usepackage[margin=25mm]{geometry}
\usepackage{fontspec}
\setmainfont{Latin Modern Roman}
\usepackage{amsmath,amssymb,mathtools,unicode-math}
\setmathfont{Latin Modern Math}
\usepackage{enumitem,needspace,xurl,hyperref}
\setlength{\parindent}{0pt}
\setlength{\parskip}{5pt}
\setlength{\emergencystretch}{3em}
\newcommand{\N}{\mathbb{N}}
\newcommand{\Z}{\mathbb{Z}}
\newcommand{\Q}{\mathbb{Q}}
\newcommand{\R}{\mathbb{R}}
\newcommand{\C}{\mathbb{C}}
\newcommand{\F}{\mathbb{F}}
\newcommand{\A}{\mathbb{A}}
\newcommand{\PP}{\mathbb{P}}
\newcommand{\E}{\mathbb{E}}
\newcommand{\eps}{\varepsilon}
\newcommand{\dd}{\,\mathrm{d}}
\begin{document}
\section*{UnsolvedMath ID: 10400101}
\textbf{Source problem code:} catalogue entry
\subsection{Definitions and mathematical statement}
For an admissible integer \(n\), let \(S_n\) be the symmetric group on \(n\) letters. For an integer parameter \(m\), the source denotes by \(S_n^m\) the quandle with underlying set \(S_n\) and operation
\[x*y=y^{-m}xy^m.\]
The requested task, attributed to J. S. Carter, is to compute the quandle groups printed in the extraction as \(H\ Q\ i\ (S_n^m)\), uniformly in the intended \(n,m,i\). The extraction does not preserve whether \(i\) is a subscript or superscript; the research summary writes \(H_Q^i\), whereas the displayed extraction does not decide homology versus cohomology. The coefficient group and allowed parameter ranges are also not supplied.

For precision about the two possible targets, under the usual integral normalized quandle-chain convention let \(C_k^Q(X)\) be the free abelian group on \(X^k\), modulo tuples having two equal adjacent entries, with
\[\begin{split}
\partial_k(x_1,\ldots,x_k)=\sum_{j=2}^k(-1)^j\bigl(& (x_1,\ldots,\widehat{x_j},\ldots,x_k)\\
 &-(x_1*x_j,\ldots,x_{j-1}*x_j,x_{j+1},\ldots,x_k)\bigr).
\end{split}\]
Here the hat means omission. Then \(H_i^Q(X;\Z)=\ker\partial_i/\operatorname{im}\partial_{i+1}\). With an abelian coefficient group \(A\), ordinary quandle cohomology instead uses \(\operatorname{Hom}_{\Z}(C_*^Q(X),A)\) and its induced coboundary. These identify conventional candidate meanings, not a choice between the source's unresolved notations or a claim that twisted coefficients are intended.

The following section supplies motivation through the quandle cocycle invariant. Given \(\alpha\in H_Q^2(X;A)\), choose a representative two-cocycle \(\phi\), with \(A\) written multiplicatively. Color oriented knot-diagram over-arcs by \(X\), equivalently use the indicated representations \(Q(K)\to X\). At the two crossing configurations in (35), assign weights \(\phi(x,y)\) and \(\phi(x,y)^{-1}\), respectively, according to the source's missing pictures. It states
\[\Phi_\alpha(K)=\sum_C\prod_\tau W(\tau,C)\in\Z[A],\]
where \(C\) ranges over colorings and \(\tau\) over crossings, and that this depends only on \(K\) and \(\alpha\). The sum is finite for \(X=S_n^m\); for unrestricted \(X\), finiteness or a completion would need specification. This following exposition is background rather than a second requested computation.
\subsection{Short English statement}
Compute the quandle homological or cohomological groups of the symmetric group equipped with conjugation by the \(m\)th power. The source extraction leaves the index position and coefficients unclear; the following cocycle-invariant discussion explains the motivation.
\subsection{Sources}
\begin{itemize}
\item Local UnsolvedMath registry, permanent ID \(10400101\), exact source code \texttt{catalogue entry}.
\item J. S. Carter, attribution in Ohtsuki, \emph{Problems on invariants of knots and 3-manifolds} (2002), Problem 5.5, PDF page 88: \url{https://msp.org/gtm/2002/04/gtm-2002-04-024s.pdf}. Recorded access 2026-07-29, pdf-font-structure extraction.
\item The following Section 5.4 attributes the cocycle invariant to internal references [79, 80]; their full bibliographic entries are not supplied.
\end{itemize}
\end{document}
